% Tapis de Quillen, with Quillen's offprint and preprint (cote 162-5, typed part): archive cote 162-5, site page 40 (PDF page 41).
% typescript. Draft compared: none.
% AI-transcribed by Claude Opus 5.5 (claude-opus-5-5, adaptive thinking, effort high), Message Batches API, batch t162_full (msgbatch_01UAh3zyWPd454AeDkhnmQ33), 2026-11-13; not yet checked by eye.
\scanpage{162-5/40}
Ces\xout{x} bicomplexes s'identifient aux complexes de chaines de longueur 2.
Donc la catégorie de (2-Cat)-groupes abéliens est équivelente à celle des
complexes de chaines de longueur 2. L'explicitation de la 2-catégorie
associée à un tel complexe de chaines est mon interprétation chérie. Il
faudrait encore expliciter la\xout{sxmorphismex} 2-Grab-catégorie des
homomorphismes (en un sens à préciser) entre deux\xout{x} telles 2-Grab-catégo-%
ries strictes, en termes d'un complexe RHom.

L'idée qui se dégage de cette approche, et que Quillen on espère va
expliciter, c'est que la \xout{notionxdexcatégoriexde} notion\xout{s} de groupe abélien
dans (n-Cat) correspond à celle de complexe de \inserted{chaines de} longueur n, en raisonnant
soit mod. isomorphisme de part et d'autre, soit mod. une notion naturelle
de n-équivalence d'un côté, et dans la catégorie dérivée de l'autre. Ceci
doit pouvoir se préciser en utilisant des Hom$^{\bullet}$ resp. des RHom$^{\bullet}$.

Si maintenant C est une catégorie avec un ``produit tensoriel'' associa
tif (et éventuellement unitaire), i.e. un bifoncteur avec des isomorphis-%
mes d'associativité (mais sans inverse ...), on peut lui associer sans
doute une Gr-catégorie universelle (tout comme à un monoïde on associe un
groupe symmétrisé), d'où des invariants, qu'on pourra désigner par $K_o(C)$
et $K_1(C)$. On suppose qu'on pourra de même trouver, pour tout n, une
Gr-n-catégorie universelle où envoyer C de façon compatible avec $\otimes$ (tout
cela dans un sens à préciser), dont les invariants $\hand{$\pi$}_n$ seraient par défini
tion les $K_n(C)$. Un cas intéressant est celui où C est une catégorie avec
produits finis (par exemple une catégorie additive) et $\otimes$ est le produit,
\xout{maisx} \over{(}{[}Mais on notera que dans le cas d'une catégorie non seulement additive
mais abélienne, il y a une notion plus adéquate du $K_o$ , qui devrait corres
pondre à une notion plus adéquate de la suite des $K_i$ , inspirée du tapis
de Roos; sans doute c'est en termes de (``vraie'') catégories triangulées

\mspage{40}
